\documentclass[10pt,a4paper]{amsart}
\usepackage[margin=19mm]{geometry}
\usepackage{amsmath,amssymb,mathtools}
\usepackage[scr=rsfs,cal=euler]{mathalfa}
\usepackage{xfrac}
\usepackage[unicode,hidelinks,bookmarks=false]{hyperref}
\newcommand{\ie}{i.e.\ }
\newcommand{\cf}{cf.\ }
\newcommand{\resp}{respectively\ }
\newcommand{\Subsection}{\subsection}
\providecommand{\nobreakdash}{\nobreak\hbox{-}}
\setlength{\emergencystretch}{2em}
\multlinegap0pt
\usepackage{amssymb}
\usepackage{xcolor}
\newtheorem{prop}{Proposition}
\DeclareMathOperator{\Fix}{Fix}
\DeclareMathOperator{\Frob}{Frob}
\newcommand{\calZ}{\mathcal{Z}}
\renewcommand{\ge}{\geqslant}
\renewcommand{\le}{\leqslant}
\title{Sharp lower bounds for shifted moments of Dedekind zeta functions}
\author{Benjamin Durkan, Nilmoni Karak, Kamalakshya Mahatab}
\date{Source: arXiv:2609.01101v1 (CC BY 4.0)}
\begin{document}
\maketitle

\noindent\emph{Source and licence.} Sharp lower bounds for shifted moments of Dedekind zeta functions. Version arXiv:2609.01101v1 (CC BY 4.0), distributed by arXiv under the Creative Commons Attribution 4.0 International licence, http://creativecommons.org/licenses/by/4.0/, as recorded in the arXiv licence metadata for this exact version and shown on the abstract page https://arxiv.org/abs/2609.01101v1. Full licence notice: \texttt{licenses/arxiv-2609.01101-CC-BY-4.0.txt}.\par\smallskip

\noindent\textit{Editorial scope.} Counted unit: the article's prop prop:square-series-B (source0.tex lines 443--465). The article's complete original proof of it is delivered with it (source0.tex lines 467--565, one \texttt{proof} environment of 99 lines). No mathematics is written, repaired or re-derived: the statement and the proof are the article's own lines, in the article's own notation, and no retained line is substituted or re-flowed.  Retained uncounted as the source-local context the proof needs: lines 136--142; lines 393--396; lines 398--402.  Nothing was omitted from any retained span, so the package carries no omission record.  No retained line contains a citation, so the wrapper carries no bibliography and no bibliography entry.  No retained line contains a figure environment, an image-inclusion command or drawing code, so no image is delivered and the package declares no asset dependency.  Editorial formatting only, in addition: an AMS \texttt{amsart} preamble that restates the article's own declarations and macros used by the retained lines, namely the environments \texttt{prop} on separate counters, together with the standard packages those lines need.  The retained environments print in this document as Proposition 1 in order of appearance, whereas the article prints the same items with the numbering of its own full document.  The complete native source of the article as distributed in the arXiv e-print for this exact version is bundled unchanged as \texttt{sources/209081/source0.tex}.
\begin{equation}\label{eq:B-definition}
B(T,\mathbf b,\mathbf a)=
\prod_{i,j=1}^r
\left|
\calZ_{ij}\left(1+\frac{1}{\log T}+i(b_i-b_j)\right)
\right|^{a_i a_j/4}.
\end{equation}

\begin{equation}\label{eq:first-prime-coefficient}
c_{\mathbf b}(p)=
\sum_{i=1}^r q_i\chi_i(\Frob_p)p^{-ib_i}.
\end{equation}

\begin{equation}\label{eq:square-euler-series}
\log\sum_{n=1}^{\infty}
\frac{|c_{\mathbf b}(n)|^2}{n^\sigma}
=\sum_p\frac{|c_{\mathbf b}(p)|^2}{p^\sigma}+O(1),
\end{equation}

\begin{prop}\label{prop:square-series-B}
There is $A_0\ge1$ such that, uniformly for
\begin{equation*}
1\le\lambda\le\tfrac12\log T
\qquad\hbox{and}\qquad |b_i|\le T/2,
\end{equation*}
one has
\begin{equation}\label{eq:H-lambda-comparison}
(2+\lambda)^{-A_0}B(T,\mathbf b,\mathbf a)
\ll H_\lambda(T)
\ll(2+\lambda)^{A_0}B(T,\mathbf b,\mathbf a),
\end{equation}
where
\begin{equation}\label{eq:H-lambda-definition}
H_\lambda(T)=\sum_{n=1}^{\infty}
\frac{|c_{\mathbf b}(n)|^2}{n^{1+\lambda/\log T}}.
\end{equation}
Also, for some $A_1\ge1$,
\begin{equation}\label{eq:B-poly}
(\log T)^{-A_1}\ll B(T,\mathbf b,\mathbf a)
\ll(\log T)^{A_1}.
\end{equation}
\end{prop}

\begin{proof}
For $p\notin S_L$, equation~\eqref{eq:first-prime-coefficient} gives
\begin{equation}\label{eq:prime-square-expansion}
|c_{\mathbf b}(p)|^2
=\frac14\sum_{i,j=1}^r
a_i a_j\chi_i(\Frob_p)\chi_j(\Frob_p)
p^{-i(b_i-b_j)}.
\end{equation}
The diagonal action of $G$ on $G/H_i\times G/H_j$ has character
$\chi_i\chi_j$.  Its orbits are indexed by $H_i\backslash G/H_j$, and
the stabiliser of $(H_i,gH_j)$ is $H_i\cap gH_jg^{-1}$.  Decomposing this
$G$-set into its orbits gives, for every $\sigma\in G$,
\begin{equation}\label{eq:mackey-character-identity}
\begin{split}
\chi_i(\sigma)\chi_j(\sigma)
&=\sum_{H_i gH_j\in H_i\backslash G/H_j}
\Bigl|\Fix\Bigl(\sigma:G/(H_i\cap gH_jg^{-1})\\
&\hspace{9em}\longrightarrow
G/(H_i\cap gH_jg^{-1})\Bigr)\Bigr|.
\end{split}
\end{equation}
For $\Re s>1$, take the logarithms defined by the absolutely convergent
Euler products.  At an unramified prime $p$, the coefficient of $p^{-s}$
in the logarithm of the factor attached to
$L^{H_i\cap gH_jg^{-1}}$ is the number of fixed points of $\Frob_p$ on
$G/(H_i\cap gH_jg^{-1})$.  It follows from
\eqref{eq:mackey-character-identity} that
\begin{equation}\label{eq:double-coset-log-euler}
\log\calZ_{ij}(s)
=\sum_{p\notin S_L}
\frac{\chi_i(\Frob_p)\chi_j(\Frob_p)}{p^s}+O(1),
\end{equation}
where the $O(1)$ term is uniform for $\Re s>1$.  Indeed, the ramified
primes form a fixed
finite set, and the terms of prime-power exponent at least two are bounded
in absolute value by a fixed multiple of
$\sum_p\sum_{m\ge2}p^{-m\Re s}$.

Set $s=1+1/\log x+iu$.  Since the character values are bounded, Mertens'
estimate and partial summation give
\begin{equation*}
\sum_{p\le x}\frac{1-p^{-1/\log x}}p
+\sum_{p>x}\frac1{p^{1+1/\log x}}\ll1.
\end{equation*}
Comparing \eqref{eq:double-coset-log-euler} with the truncated prime sum
and taking real parts therefore gives, uniformly for real $u$ and $x\ge3$,
\begin{equation}\label{eq:double-coset-prime-sum}
\sum_{\substack{p\le x\\p\notin S_L}}
\frac{\chi_i(\Frob_p)\chi_j(\Frob_p)
\cos(u\log p)}{p}
=\log\left|\calZ_{ij}
\left(1+\frac{1}{\log x}+iu\right)\right|+O(1).
\end{equation}
The terms indexed by $(i,j)$ and $(j,i)$ in
\eqref{eq:prime-square-expansion} are complex conjugates, so their
imaginary parts cancel when the ordered pairs are summed.
Equations~\eqref{eq:prime-square-expansion} and
\eqref{eq:double-coset-prime-sum}, summed over the ordered pairs $(i,j)$,
give
\begin{equation}\label{eq:prime-sum-B}
\sum_{p\le T}\frac{|c_{\mathbf b}(p)|^2}{p}
=\log B(T,\mathbf b,\mathbf a)+O(1).
\end{equation}

Let $\delta=\lambda/\log T$.  Since
$|c_{\mathbf b}(p)|\ll1$, Mertens' estimate and partial summation give
\begin{align}
&\sum_{p\le T}\frac{1-p^{-\delta}}p
+\sum_{p>T}\frac1{p^{1+\delta}}
\ll \log(2+\lambda).                                    \label{eq:smoothing-cost}
\end{align}
For example, the first sum is bounded by a constant multiple of
\begin{equation*}
1+\int_{1/\log T}^{1}\frac{1-e^{-\lambda v}}v\,dv,
\end{equation*}
and the second by a constant multiple of
$1+\int_1^\infty e^{-\lambda v}v^{-1}\,dv$.
Combining \eqref{eq:square-euler-series}, \eqref{eq:prime-sum-B}, and
\eqref{eq:smoothing-cost} yields
\begin{equation*}
\log H_\lambda(T)
=\log B(T,\mathbf b,\mathbf a)
+O(\log(2+\lambda)).
\end{equation*}
This proves \eqref{eq:H-lambda-comparison}.

Every Dedekind zeta function in \eqref{eq:B-definition}, evaluated at
$1+1/\log T+iu$, has modulus at most its value at
$1+1/\log T$, which is $O(\log T)$.  If $F$ is one of the fixed fields
occurring in that definition and $\sigma=1+1/\log T$, then its Euler
product also gives
\begin{equation*}
|\zeta_F(\sigma+iu)|^{-1}
\le \prod_{\mathfrak p}\left(1+(N\mathfrak p)^{-\sigma}\right)
=\frac{\zeta_F(\sigma)}{\zeta_F(2\sigma)}
\ll\log T.
\end{equation*}
There are only finitely many factors.  This proves \eqref{eq:B-poly}.
\end{proof}

\end{document}
